\chapter{What survives when the model changes}\label{ch:robustness}

\begin{chapterplan}
Propositions A.3 to A.6 change one feature of the base model at a time. Each proof follows Proposition~2: orders
first, then posteriors, entry and prices.
\end{chapterplan}

%==================================================================
\section{Preparation cost laws}\label{sec:robustness-costs}

The cost law mixes a low and a high component, \(\Phi=\rho\Phi^0+(1-\rho)\Phi^1\), with
\(\operatorname{supp}\Phi^0\subset[0,\bar c_0]\) and \(\operatorname{supp}\Phi^1\subset[c_-,c_+]\). Write \(B_i=B_{F_i}\),
\(\Delta_i=\DT(F_i)\), \(m=m(\pi)\), \(M=M(\pi)\) and \(\kappa=\kappa(\pi)=\min\{m,1-M\}\).

\begin{paperresult}{Proposition A.3 (Preparation cost laws)}
Let the classes satisfy the class ranking, and let \(F_1\) strengthen \(F_0\). The cost law is
\(\Phi=\rho\Phi^0+(1-\rho)\Phi^1\) as above. Suppose
\[
0<\bar c_0<B_1(m),\qquad B_0(\pi)<c_-<c_+<B_1(M),\qquad \Delta_0<k<\rho\kappa\Delta_1.
\]
\begin{enumerate}[label=(\roman*),leftmargin=2em]
\item At \(F_0\), the unique orders are zero, and entry is \(\rho\).
\item At \(F_1\), the unique orders are full, and only a purchase recruits the high component. Entry is \(\rho+(1-\rho)(b+\pi\lambda)\). High-class ownership rises if and only if \([\rho+(1-\rho)(b+\lambda)]\,W_H(F_1)>\rho\,W_H(F_0)\).
\end{enumerate}
\end{paperresult}

\emph{Proof.} The new step is the entry floor. A price belief \(\bar\mu\) at a nonzero order is at least \(m\), so
\(B_i(\bar\mu)\ge B_i(m)\ge B_1(m)>\bar c_0\): every low cost prepares, and \(e(\pm1)\ge\rho\). The trading bounds give
zero orders at \(F_0\) and full orders at \(F_1\). Only a purchase covers the high costs, \(B_1(M)>c_+\), so entry
is \((\rho,\rho,1)\). A profit margin over every low cost replaces the zero cost.

%==================================================================
\section{Value classes and priors}\label{sec:robustness-classes}

\begin{paperresult}{Proposition A.4 (Value classes and priors)}
Fix classes that satisfy the class ranking with \(\E_H[\theta]>\E_L[\theta]\), a reserve \(p>0\), and any \(\pi,\rho\in(0,1)\) and slot \((\lambda,\beta)\). For uniform incumbents, a nonempty open set of strengths and costs \((r_0,r_1,r_2,c_H,k)\) satisfies (H1)--(H3) and the ownership condition of Proposition~2. On this set every conclusion of Proposition~2 holds.
\end{paperresult}

\emph{Proof.} Let \(D=G_L-G_H\), \(\Gamma(x)=\int_p^xD(y)\,dy\), \(\bar\Gamma=\E_H[\theta]-\E_L[\theta]\) and
\(\theta_*=\inf\{x\ge p:\Gamma(x)>0\}\). Against \(U[0,r]\),
\begin{align*}
\DT(r)&=\frac1r\int_p^r\Gamma(y)\,dy\le\frac{(r-\theta_*)\,\bar\Gamma}{r}\to0\quad(r\downarrow\theta_*) \why{integration by parts; \(\Gamma=0\) below \(\theta_*\)}\\
g_H(r)-g_L(r)&=\bar\Gamma-\DT(r)\ge p\,\bar\Gamma/r>0 \why{class-gap identity}
\end{align*}
So pick \(r_0<r_1\) just above \(\theta_*\), \(c_H\) and \(k\) inside the two open windows, and \(r_2\) large.

\begin{watchout}
Some costs work for every prior, but the example's costs do not work for every prior.
\end{watchout}

%==================================================================
\section{Many classes under separation}\label{sec:robustness-many}

With many classes the posterior is a vector, but the price is one number. Nonzero orders give posterior vectors in \(\mathcal M_\pi=\{\nu:\nu_j/\pi_j\le(1+u)\,\nu_l/\pi_l\ \text{for all }j,l\}\). Split the classes into a high
group \(\mathcal H\) and a low group \(\mathcal L\), let \(\bar\pi=\sum_{j\in\mathcal H}\pi_j\), and let \(\bar g_{\mathcal
H},\bar g_{\mathcal L}\) be the prior-weighted group means of \(g_{j,1}\).

\begin{paperresult}{Proposition A.5 (Many classes under separation)}
Let the incumbent strengthen from \(F_0\) to \(F_1\). Suppose
\[
\max_jt_{j,0}-\min_jt_{j,0}<k,\qquad \sum_j\pi_jg_{j,0}<c_H,
\]
\[
\min_{j\in\mathcal H}t_{j,1}-\max_{\nu\in\mathcal M_\pi}\nu\cdot t_1>\frac k\rho,\qquad
\min_{\nu\in\mathcal M_\pi}\nu\cdot t_1-\max_{j\in\mathcal L}t_{j,1}>\frac k\rho,
\]
and \(\bar g_{\mathcal H}>\bar g_{\mathcal L}\) and \(\bar g_{\mathcal L}+M(\bar\pi)(\bar g_{\mathcal H}-\bar g_{\mathcal L})>c_H\).
Then at \(F_0\) the unique orders are zero, and entry is \(\rho\). At \(F_1\) the unique orders are a purchase in every high-group class and a sale in every low-group class, only a purchase recruits, and entry is \(\rho+(1-\rho)(b+\bar\pi\lambda)>\rho\). These conclusions hold under either challenger rule.
\end{paperresult}

\emph{Proof.} Class \(j\) earns \(U_j(q)=q\,e(q)[t_j-\mu(q)\cdot t]-k|q|\), with \(e\in[\rho,1]\). At \(F_0\) every
nonzero order loses. At \(F_1\), separation makes the purchase strictly best in \(\mathcal H\) and the sale in
\(\mathcal L\). Each group then shares one order law, so the group posterior \(\mu_{\mathcal H}\) is the two-class
posterior at the prior \(\bar\pi\), and \(B(\mu)=\bar g_{\mathcal L}+\mu_{\mathcal H}(\bar g_{\mathcal H}-\bar g_{\mathcal
L})\). The price reveals \(\mu_{\mathcal H}\). Without separation, one price can carry two vectors.

%==================================================================
\section{Complementary private signals}\label{sec:robustness-signals}

Proposition A.6 gives the investor a signal \(S_I\) and the challenger a signal \(S_C\), with accuracies
\(\alpha_I,\alpha_C\in(\tfrac12,1)\). A purchase gives the public posterior \(M^I=(1-\alpha_I)+(2\alpha_I-1)M(\tfrac12)\), and \(\mu^C_+(\mu)\) is the challenger's belief
after \(S_C=+\). The paper's conditions are
\[
B_{F_0}(\alpha_C)<c_H<B_{F_1}(\mu^C_+(M^I)),\qquad
(2\alpha_I-1)\bar\Delta_P(F_0)<k<(2\alpha_I-1)\rho\,m(\tfrac12)\DT(F_1),
\]
with \(\bar\Delta_P(F)=\DT(F)+(1-\rho)(2\alpha_C-1)[t_H(F)-t_0(F)]\). Then, under all four tie rules, the investor
abstains at \(F_0\), with entry \(\rho\); at \(F_1\) it trades on its signal, and entry exceeds \(\rho\).
The price slope \(\Delta_P=e_L\DT+(e_H-e_L)w_H\in[\rho\DT,\bar\Delta_P]\) bounds the investor's returns. In the paper's examples both sources are needed: a favorable \(S_C\) alone does not cover \(c_H\).

\begin{chaptersummary}
The core of Proposition~2 survives each change. A low cost range replaces the zero cost when profit at the lowest
feasible belief covers it. Every pair of ranked classes with different means and every prior admit some costs and strengths that work.
Separation lets one price reveal a posterior vector, and private signals complement the price.
\end{chaptersummary}
